\begin{algorithm}\footnotesize
\caption{LazyBoE}\label{algo:lazyboe}
  \KwIn{Start State $q_0$, Goal Region $\mathcal{Q}_{goal}$, Edge Bundle $\mathcal{E}$, Neighborhood Radius $\theta$}
  \KwOut{Path $\pi$}
  $\pi \leftarrow \emptyset$\\
  $\mathcal{V} \leftarrow \{q_0\}$\\
  \While{$\mathcal{V} \neq \emptyset$} {
    \tcp{equivalent to SELECT(.) in \cite{shome2021asymptotically}}
    $v \leftarrow \text{PICK\_NODE}(\mathcal{V})$\\
    \tcp{update $\pi$ to lower cost path}
    \If{$v \in \mathcal{Q}_{goal}$ \textbf{and} $!\text{IS\_LAZY}(\text{EDGE\_TO}(v))$} {
        \If{cost$(\pi) > $ cost$(\text{PATH\_TO}(v))$} {
            $\pi \leftarrow \text{PATH\_TO}(v)$
        }
    }
    \If{IS\_LAZY(EDGE\_TO(v))} {
        $\mathcal{E}_{\text{neighbor}} \leftarrow \text{RADIAL\_NN}(v, \theta / 2, \mathcal{E})$
    }
    \Else {
        $\mathcal{E}_{\text{neighbor}} \leftarrow \text{RADIAL\_NN}(v, \theta, \mathcal{E})$
    }
    %\\
    $\mathcal{V}_{\text{next}} \leftarrow \emptyset$\\
    \ForEach{$e \in \mathcal{E}_{\text{neighbor}}$} {
        $P_{\text{select}} \leftarrow (1 - e.P_{\text{collision}}) * e.P_{\text{lazy\_prop}}$\\
        $\text{apply\_lazy\_prop} \leftarrow rand() < P_{\text{select}}$\\
        \If{apply\_lazy\_prop} {
            $\mathcal{V}_{\text{next}} \leftarrow \mathcal{V}_{\text{next}} \cup e.q_f$\\
        }
        \ElseIf{!apply\_lazy\_prop \textbf{or} $\text{!IS\_LAZY(EDGE\_TO(v))}$ \textbf{or} (IS\_LAZY(EDGE\_TO(v)) \textbf{and} SHOULD\_TERMINATE\_LAZY())} {
            \tcp{simulate}
            $e_{new} \leftarrow f(v, e.u, e.\Delta t)$\\
            \If{$\text{COLLISION\_FREE}(e_{new})$ \\
   \textbf{and} $\text{WITHIN\_LIMITS}(e_{new})$ \\
   \textbf{and} $\text{IS\_CONTINUOUS}(e_{new})$} {
                $\mathcal{V}_{\text{next}} \leftarrow \mathcal{V}_{\text{next}} \cup e_{new}.q_f$
            }
        }
    }
    $\mathcal{V} \leftarrow \mathcal{V} \cup \mathcal{V}_{\text{next}}$

  }
  \textbf{return} $\pi$ 
\end{algorithm}


\subsection{Edge Perturbation for Lazy Propagation and Collision-Checking}\label{sec:methods-perturbation}

The task of generating an edge bundle that uniformly samples the state and control spaces of a robot is a computationally intensive process, particularly in the case of a high-DoF robot and increased environment complexity.
A substantial part of this computational workload is incurred again during the planning stage, where edges emerging from $\theta$-regions around planning tree nodes undergo simulation and collision checks \cite{shome2021asymptotically}. This $\theta$ is dependent on state-space dimensions and denotes the maximum level of similarity in the state space, thereby bounding the set of valid controls.
Additionally, the greedy, heuristic-biased search strategy used by the BoE planner has the potential to get stuck in local minima, thereby necessitating a method for quickly evaluating these greedy paths without simulating them entirely. This motivates the need for lazy approaches to simulation and collision-checking in the kinodynamic setting.
More specifically, we need to address the following two questions for each edge $e_i \in \mathcal{E}$:
\begin{enumerate}
    \item Can $e_i$ be lazily propagated instead of undergoing a full simulation?
    \item What is the probability that a lazily propagated edge will result in a collision?
\end{enumerate}

We approach the answers to these questions empirically by conducting a perturbation analysis on each edge in the bundle. A disturbance $\Delta q$ in the range $(0, \theta]$ is added to the start state for each edge. The resultant end state for an $e_i$ is computed as $\hat{q}_f^i = f(q_0^i + \Delta q, \boldsymbol{\dot{q}}^i)$. In the limit that the number of perturbations for each edge ($p$) approaches infinity, we are able to accurately estimate both the maximum error in the end state, $\Delta q_f^i = \text{max}(\|e.\boldsymbol{q}_f - \hat{\boldsymbol{q}}_f)\|_2$), and the potential collision likelihood.
To address the first question, an edge viable for lazy propagation is characterized as one that confines the end-state error within the region radius $\theta$, so that viable edges for the next level of propagations can be picked from this fixed $\theta$-neighborhood.
However, this criterion is subjected to a more stringent requirement, as depicted in \cref{fig:theta_over_2}, i.e., the end-state error for $e_i$ must be confined to $\theta / 2$ if $e_i$ undergoes lazy propagation to ensure the worst-case error between its real simulation and the following connection to the lazy edge is $\theta$.
Put differently, the probability that an edge can be lazily propagated $P_{\text{lazy\_prop}}^i = Pr(\Delta q_f^i < \theta / 2)$ is the ratio of the number of valid (collision-free and dynamically-feasible) perturbations for which the maximum end state error is $\theta / 2$ to the total number of valid perturbations, $p_{valid}$ (\cref{fig:pert_analysis}a).
Similarly, in response to the second question, we can quantify the probability that an edge will collide when lazily propagated, i.e., $P_{\text{collision}}^i = Pr(e_i \text{ collides})$, as the ratio of the number of edges that collide to the total number of perturbations $p$ (\cref{fig:pert_analysis}b). This is in contrast to prior work which only considers lazy collision checking in the geometric planning domain and takes a deterministic approach, iteratively discarding edges that are not collision-free \cite{kavraki2000path,hauser2015lazy}. The neighborhood-based control propagation in our work allows for a more nuanced and stochastic approach to lazy collision checking by allowing edges to be re-explored, potentially uncovering paths that would be prematurely eliminated by deterministic methods. In this work, we primarily deal with self-collisions and environmental collisions between the robot and its surrounding environment.
These two probability values for each edge extend the definition of an edge bundle $\mathcal{E}$ to the following:
\begin{center}
$\mathcal{E} = \{(q_0, \boldsymbol{\dot{q}}, \Delta t, q_f, P_{\text{lazy\_prop}}, P_{\text{collision}}) \mid q_0, q_f \in \mathcal{Q}, \hspace{3pt} \boldsymbol{\dot{q}} \in \mathcal{U}\}$
\end{center}

These two metrics, associated with each edge, are used to guide the planner by deciding when to apply lazy methods for faster planning times.
